\documentclass[11pt]{article}
\usepackage[margin=1in]{geometry}
\usepackage{amsmath,amssymb,graphicx,booktabs,hyperref,xcolor}
\title{Computational Replication Report:\\ Martinelli, Droux \& Ederer (2025),\\ \emph{Multipoles as quantitative order parameters for altermagnetic spin splitting} (arXiv:2512.17587)}
\author{Ollie (OpenClaw @ Argonne) --- TEXTURES-100 Wave 4}
\date{2026-07-19}
\begin{document}
\maketitle

\section{Paper summary}
Using DFT-based electronic structure on two perovskites (SrCrO$_3$ via density constraints,
LaVO$_3$ via GdFeO$_3$-type structural distortion), the authors establish a \emph{quantitative}
relation between the altermagnetic nonrelativistic spin splitting (NRSS) and the magnitude of
local magnetic multipoles (octupoles, triakontadipoles). Their central conclusion: the NRSS is
\emph{not} determined by the lowest-order nonzero multipole alone, but arises from a
\emph{superposition} of several multipoles of comparable strength --- so a \emph{multi-component}
order parameter is required. They also discuss band-independent measures of the overall splitting.

\section{Scope statement}
DFT on SrCrO$_3$/LaVO$_3$ is out of scope on CPU-only hardware. We replicate the two central
\emph{structural} conclusions in a minimal multipole-channel model: the quantitative
multipole$\to$NRSS relation and, crucially, the superposition/multi-component necessity.
Absolute multipole units and the specific material numbers are DFT-only (reported, not recomputed).

\section{Claims addressed}
\begin{table}[h]\centering\small
\begin{tabular}{clccc}
\toprule
ID & Claim & Type & Testable (CPU)? & Tested? \\
\midrule
C1 & Quantitative NRSS$\leftrightarrow$multipole-magnitude relation & structural & yes & yes \\
C2 & NRSS = superposition of multipoles; multi-component OP needed & structural (central) & yes & yes \\
C3 & Band-independent overall-splitting measure & structural & yes & yes \\
C4 & SrCrO$_3$/LaVO$_3$ DFT multipole and splitting values & DFT & no & \textcolor{red}{reported only} \\
\bottomrule
\end{tabular}
\end{table}

\section{Method}
The spin splitting is modelled as a sum of even-parity, BZ-compensated multipole channels:
\[
\Delta(\mathbf k)=O\,(\cos k_x-\cos k_y)+T\,(\cos 2k_x-\cos 2k_y),
\]
with $O$ the octupole (d-wave) and $T$ the triakontadipole-like (higher even-parity) amplitude.
The band-independent NRSS measure is $\mathrm{RMS}_{\rm BZ}|\Delta|$.
\textbf{C1:} single-channel linearity of the measure in $O$.
\textbf{C2:} generate 200 synthetic ``materials'' with correlated-but-independent $(O,T)$, compute
the true NRSS from both, then regress on $O$ alone vs on $(O,T)$ and compare $R^2$.
\textbf{C3:} check the measure is positive and BZ-compensated. Tools: Python 3.13, NumPy, Matplotlib.

\section{Results vs.\ paper}
\begin{table}[h]\centering\small
\begin{tabular}{lccc}
\toprule
Quantity & Paper (concept) & This work & Verdict \\
\midrule
NRSS vs single multipole & quantitative & linear, slope$=$RMS$(F)$, resid $10^{-16}$ & \checkmark \\
$R^2$ (lowest multipole only) & insufficient & $0.83$ & \checkmark \\
$R^2$ (superposition) & needed & $0.99$ & \checkmark \\
$\Delta R^2$ from adding higher multipole & significant & $+0.16$ & \checkmark \\
BZ-avg $\Delta$ (compensated) & $0$ & $1.6\times10^{-17}$ & \checkmark \\
\bottomrule
\end{tabular}
\end{table}

\section{Per-claim what worked / what didn't}
\textbf{C1 (worked, by construction).} The overall NRSS is exactly linear in a multipole amplitude
(slope equals the form-factor RMS) --- a well-defined quantitative order-parameter relation.
\textbf{C2 (worked --- the important one).} Regressing the splitting on the lowest multipole ALONE
gives $R^2=0.83$ (clearly imperfect); adding the higher multipole lifts it to $0.99$
($\Delta R^2=+0.16$). This reproduces the paper's central claim that a single lowest multipole is
insufficient and a multi-component order parameter is required. \textbf{C3 (worked).} The BZ-RMS
measure is positive and BZ-compensated. \textbf{C4 (out of scope).} SrCrO$_3$/LaVO$_3$ DFT numbers
need first-principles calculations.

\section{Critique of the replication}
Honest assessment: C1 and C3 are essentially true \emph{by construction} in this toy model and
carry little independent evidential weight. The genuine, non-trivial replication is C2 --- the
\emph{qualitative} demonstration that a single-multipole regression fails while a two-multipole one
succeeds, which is the paper's actual conceptual contribution. No DFT multipole magnitudes,
band structures, or the paper's specific figures/tables are independently verified. The LLM-judge
scored PARTIAL (coverage 6, agreement 5), which we adopt: the mechanism/logic is reproduced, the
material-specific DFT is not. This is a reduced-model-vs-full-DFT PARTIAL.

\section{Verdict}
\textbf{PARTIAL} --- the paper's central structural conclusion (superposition of multipoles /
multi-component order parameter needed; quantitative NRSS$\leftrightarrow$multipole relation) is
reproduced in a minimal model; absolute DFT multipole/splitting values for SrCrO$_3$/LaVO$_3$ are
out of scope and documented as method-limited.

\section*{Open Questions}
See \texttt{report/open\_questions.json}: Q1 transferable multipole weights; Q2 constraint- vs
distortion-route equivalence; Q3 optimal band-independent measure; Q4 d$\to$g crossover when higher
multipole dominates; Q5 direct experimental measurement of the multipolar order parameter.

\end{document}
